\section{Conclusion}

We presented a large-scale study of multi-horizon regional aircraft trajectory forecasting on ADRR-derived open aviation data. Starting from a hierarchical H3 GRU baseline, we introduced a horizon-query Transformer decoder and extended it with an explicit multiscale temporal branch. The final model improves top-k accuracy, NLL, Brier score, and calibration relative to both the GRU baseline and a single-scale Transformer ablation, while also producing smaller candidate region sets. A classical Markov baseline is far less competitive, which confirms that the task requires richer temporal representation learning than simple empirical transition modeling can provide.

The H3 refinement study further shows that candidate recall and calibration should be reported together. A model that expands candidate regions aggressively can appear stronger under recall-only analysis while becoming less trustworthy probabilistically. We therefore argue that large-scale regional aircraft forecasting should be evaluated as a joint problem of accuracy, uncertainty quality, and spatial efficiency.

Overall, the results support hierarchical H3 forecasting with multiscale horizon-query temporal encoding as a strong baseline for long-horizon regional aviation prediction on open trajectory archives. They also provide a stable foundation for the next methodological step: explicitly uncertainty-aware forecasting under missing and irregular flight observations.
